\documentclass[11pt,a4paper]{article}
\usepackage[utf8]{inputenc}
\usepackage[T1]{fontenc}
\usepackage[brazil]{babel}
\usepackage[a4paper,margin=2.5cm]{geometry}
\usepackage{amsmath,amssymb}
\usepackage{enumitem}
\usepackage{xcolor}
\usepackage{hyperref}

\hypersetup{
  colorlinks=true,
  linkcolor=blue!60!black,
  urlcolor=blue!60!black
}

\usepackage{listings}
\lstset{
  basicstyle=\ttfamily\small,
  columns=fullflexible,
  keepspaces=true,
  numbers=left,
  numberstyle=\tiny\color{gray},
  frame=single,
  breaklines=true,
  language=C++
}

\title{\textbf{Roteiro para os Vídeos Individuais}\\[0.5em]
\large Rumo à Liga Pokémon --- Algoritmos em Grafos 2026.1}
\author{Gildo Alves\,\textbullet\,João Landin\,\textbullet\,Francisco Almir\\[0.3em]
\small Universidade Federal do Cariri (UFCA) --- Prof.\ Carlos Vinicius G. C. Lima}
\date{Agosto de 2026}

\begin{document}

\maketitle
\vspace{-1.5em}
\begin{center}
\rule{\textwidth}{0.4pt}
\end{center}

\begin{abstract}
\noindent Este documento é um \textbf{roteiro interno} para a gravação dos vídeos de defesa.
Cada seção indica, para cada integrante, \textbf{o que explicar}, a \textbf{ordem da apresentação}
e \textbf{quais funções abrir no código} para demonstrar a implementação ``no metal''.
Serve apenas como guia de gravação --- não é um relatório para o professor.
\end{abstract}

\section{Dicas Gerais de Gravação}

\begin{itemize}[leftmargin=2em,itemsep=2pt]
  \item Grave a sua tela mostrando o editor de código e explique lendo as funções da sua seção.
  \item Mostre \texttt{make} compilando sem erros e rode o jogo (\texttt{make run}).
  \item Foque na \textbf{sua} parte: a avaliação é individual.
  \item Cada integrante deve destacar a sua \textbf{operação sobre o grafo}: \textbf{Gildo} = Dijkstra, \textbf{João} = BFS/DFS, \textbf{Almir} = MST (Kruskal/Prim).
  \item Seja objetivo e não leia o código inteiro --- destaque a lógica essencial.
\end{itemize}

\section{Gildo Alves --- Dados, Parsing e Grafos}

\subsection{O que explicar}
Como o mundo é carregado do arquivo de texto e como os \textbf{algoritmos de grafos}
são implementados manualmente, \textbf{sem bibliotecas externas} (regra do projeto).

\subsection{Ordem da apresentação}
\begin{enumerate}[leftmargin=2em,itemsep=2pt]
  \item Estruturas de dados que definem o universo do jogo (tipos, pokémon, treinador, estado).
  \item Leitura do mapa e inicialização das entidades na memória.
  \item Algoritmos de grafos: caminho mínimo, conectividade e árvore geradora mínima.
  \item Fechamento: compilar e rodar o jogo.
\end{enumerate}

\subsection{Funções para abrir no código}
\begin{description}[leftmargin=6em,style=nextline]
  \item[\texttt{catalogo\_especies} / \texttt{multiplicador\_tipo}] constantes e tabelas em \texttt{Types.hpp}.
  \item[\texttt{carregarMapa}] leitura do \texttt{graph.txt}, montagem da lista de adjacência e sorteio das entidades.
  \item[\texttt{dijkstra}] Dijkstra com \textbf{min-heap manual} (sua operação de grafo).
  \item[\texttt{dijkstraCaminho}] menor caminho de origem a destino (base para o tempo limite $T$).
\end{description}

\subsection{Destaques técnicos}
\begin{itemize}[leftmargin=2em,itemsep=2pt]
  \item Sua \textbf{operação sobre o grafo}: caminho mínimo via \textbf{Dijkstra} (sem bibliotecas).
  \item Implementação da \textbf{fila de prioridade manual} (sem \texttt{std::priority\_queue} de grafo).
  \item Cálculo do tempo limite $T$ a partir dos pesos das arestas.
\end{itemize}

\section{João Landin --- Motor de Estados, Movimento e Tempo (+ BFS/DFS)}

\subsection{O que explicar}
Como o jogador se move no grafo, como o \textbf{tempo global} avança e como o
\textbf{estado} dos pokémon evolui conforme a distância percorrida. Sua
\textbf{operação sobre o grafo} é a validação da conectividade com \textbf{BFS} e \textbf{DFS}.

\subsection{Ordem da apresentação}
\begin{enumerate}[leftmargin=2em,itemsep=2pt]
  \item Movimentação entre vértices adjacentes e consumo de tempo.
  \item Passagem do tempo: choco de ovos, recuperação e timers de indisponibilidade.
  \item Zonas seguras (sem batalha) e gestão da equipe.
  \item Regras de inscrição no estádio (prazo e insígnias).
\end{enumerate}

\subsection{Funções para abrir no código}
\begin{description}[leftmargin=6em,style=nextline]
  \item[\texttt{moverTreinador}] deslocamento e consumo do tempo $w(e)$ da aresta.
  \item[\texttt{AvancarTempo}] relógio global e efeitos da distância percorrida.
  \item[\texttt{andandoChocarOvos}] incubação (ovo choca aos 100 unidades).
  \item[\texttt{andandoRecuperarHP}] consciente \textbf{não} recupera passivamente (função vazia de propósito).
  \item[\texttt{andandoIndisponivelReducao}] decremento do tempo de internação/indisponível.
  \item[\texttt{andandoGanharXP}] +1 XP a cada 100 unidades percorridas.
  \item[\texttt{andandoReporSelvagens}] selvagens que fugiram reaparecem quando o timer zera.
  \item[\texttt{verificar\_zonasegura}] proíbe batalha em laboratório e PMC.
  \item[\texttt{podeBatalhar}] exige 3 ou mais pokémon conscientes.
  \item[\texttt{podeInscrever} / \texttt{prazoExpirado}] inscrição no estádio com 8 insígnias antes do prazo $T$.
  \item[\texttt{podePegarOvo} / \texttt{usarErva} / \texttt{enviarParaCarvalho}] itens e gestão da equipe (máx.\ 6).
  \item[\texttt{criarJogador}] criação do treinador principal no laboratório.
  \item[\texttt{bfs} / \texttt{dfs}] busca em largura e em profundidade (sua operação de grafo).
\end{description}

\subsection{Destaques técnicos}
\begin{itemize}[leftmargin=2em,itemsep=2pt]
  \item Sua \textbf{operação sobre o grafo}: \textbf{BFS} e \textbf{DFS} para validar a conectividade e rotas de fuga.
  \item Relação entre \textbf{peso da aresta} e avanço do \textbf{relógio global}.
  \item Distinção entre reconciliações: inconsciente volta \texttt{Consciente}; \texttt{No\_PMC} exige Centro Pokémon.
  \item Regra: consciente \textbf{não} se cura apenas por caminhar.
\end{itemize}

\section{Francisco Almir --- Motor de Combate e Progressão (+ MST)}\label{sec:almir}

\subsection{O que explicar}
A matemática do combate em turnos, o cálculo de dano, a vantagem elemental,
os acertos críticos/esquivas e o sistema de evolução por XP. Sua
\textbf{operação sobre o grafo} é a \textbf{Árvore Geradora Mínima} (MST) para
análise da densidade da malha.

\subsection{Ordem da apresentação}
\begin{enumerate}[leftmargin=2em,itemsep=2pt]
  \item Pré-requisito para batalhar (equipe consciente).
  \item Cálculo de dano por turno e multiplicador de tipo.
  \item Acerto crítico, esquiva e derrubada do oponente.
  \item Ganho de XP, atributos e evolução.
  \item Batalha contra treinador e captura de pokémon selvagem.
\end{enumerate}

\subsection{Funções para abrir no código}
\begin{description}[leftmargin=6em,style=nextline]
  \item[\texttt{multiplicadorTipo}] matriz de vantagem elemental (Água, Fogo, Planta, etc.).
  \item[\texttt{aplicarDano}] aplica dano e gerencia a troca de status (inconsciente).
  \item[\texttt{deixarInconsciente}] derruba o pokémon com tempo aleatório em $[10,50]$.
  \item[\texttt{ganharXP}] concede XP conforme o resultado da batalha.
  \item[\texttt{tentarEvoluir}] evolução ao atingir XP suficiente (fase+1 e bônus de AP/DP).
  \item[\texttt{batalharTreinador}] combate em turnos entre treinadores (insígnia do ginásio).
  \item[\texttt{batalharSelvagem}] batalha contra selvagem e fase de captura.
  \item[\texttt{mstKruskal}] MST de Kruskal com union-find manual (sua operação de grafo).
  \item[\texttt{mstPrim}] MST de Prim com min-heap manual.
\end{description}

\subsection{Destaques técnicos}
\begin{itemize}[leftmargin=2em,itemsep=2pt]
  \item Sua \textbf{operação sobre o grafo}: \textbf{MST (Kruskal/Prim)} para mensurar a densidade da malha.
  \item $\text{Dano} = \max(0,\ AP_\text{atk} - DP_\text{def})$ multiplicado pelo tipo.
  \item Crítico ($2\times$) e esquiva (dano nulo) variam com a \textbf{diferença de XP}.
  \item Evolução: XP alto $\Rightarrow$ nova fase com mais ataque e defesa.
  \item Captura só de selvagem derrubado; caso contrário, ele foge e reaparece depois.
\end{itemize}

\section{Suporte Compartilhado (opcional no vídeo)}

\subsection{Funções para abrir no código}
\begin{description}[leftmargin=6em,style=nextline]
  \item[\texttt{aleatorio} / \texttt{chance} / \texttt{chance\_float}] geração de eventos probabilísticos (\texttt{RNG.cpp}).
  \item[\texttt{menuPrincipal} / \texttt{escolherAcaoBatalha} / \texttt{escolherCaptura}] menus da interface (\texttt{GUI.cpp}).
  \item[\texttt{barraHP} / \texttt{mostrarStatus} / \texttt{mostrarVizinhos}] exibição de dados (\texttt{GUI.cpp}).
\end{description}

\section{Checklist Final}
\begin{itemize}[leftmargin=2em,itemsep=2pt]
  \item \texttt{make clean \&\& make} compila sem erros de \texttt{-Wall}.
  \item Rodar \texttt{make run} (ou \texttt{./pokemon\_rpg data/graph.txt}).
  \item Cada integrante explica \textbf{apenas a sua seção}.
\end{itemize}

\end{document}
